\section{Actividades Adicionales}

Para profundizar más en la creación y edición de vídeo con DaVinci Resolve o con otra herramienta similar, a continuación se 
presentan una serie de enunciados a desarrollar de forma autodidacta y con los concoimientos consolidados sobre apartados previos de 
la documentación.

\subsection{Ejercicio 1: Carrusel sobre clips de vídeo}
Sobre un vídeo con una gran consecución de clips, se debe crear un efecto carrusel donde se pueda apreciar en formato 3D 
el paso de los clips. Es recomendable suavizar el movimiento del carrusel para obtener un resultado más fluido y amigable al 
espectador.

Material necesario:
\begin{itemize}
    \item Vídeo que contenga varios clips en secuencia.
    \item DaVinci Resolve o Plantillas de Canva video.
    \item Vídeo explicativo de la solución~\cite{carrusel}. 
\end{itemize}

\subsection{Ejercicio 2: Transición mediante desenfoque y enfoque}
Sobre dos clips de vídeo consecutivos, se debe crear una transición personalizada basada en el desenfoque progresivo del primer clip y el enfoque gradual del segundo. 
El objetivo es conseguir una transición suave que simule un cambio de plano natural. 
Se recomienda ajustar los niveles de desenfoque y sincronizar correctamente los \textit{keyframes} para evitar cortes bruscos.

Material necesario:
\begin{itemize}
    \item Dos clips de vídeo con continuidad visual.
    \item DaVinci Resolve.
    \item Vídeo explicativo de la solución~\cite{transicion-enfoque}. 
\end{itemize}

\subsection{Ejercicio 3: Construcción de una secuencia narrativa}
A partir de varios clips de vídeo sin relación aparente, se debe construir una secuencia que transmita una idea o historia coherente. 
El montaje deberá apoyarse en el uso de cortes, ritmo y selección de planos, evitando el uso excesivo de efectos. 
Se recomienda prestar especial atención a la continuidad visual y temporal para lograr un resultado comprensible para el espectador.

Material necesario:
\begin{itemize}
    \item Conjunto de clips de vídeo variados.
    \item DaVinci Resolve.
    \item Recursos explicados en el presente documento. 
\end{itemize}
